% Requires in preamble:
%   \usepackage{booktabs}
%   \usepackage[table]{xcolor}
%   \usepackage{threeparttable}

\definecolor{rowgreen}{HTML}{C8E6C9}
\definecolor{roworange}{HTML}{FFE0B2}

\begin{table}[htbp]
\centering
\caption{Overview of the seven experimental conditions spanning three axes:
         preprocessing (raw vs.\ preprocessed), label definition
         (paper vs.\ corrected), and downsampling strategy
         (none, pre-split, post-split).}
\label{tab:conditions}
\begin{threeparttable}
  \begin{tabular}{clll}
    \toprule
    \textbf{Condition} & \textbf{Preprocessing} & \textbf{Labels} & \textbf{Downsampling} \\
    \midrule
    1 & Raw          & Paper     & None       \\
    2 & Raw          & Corrected & None       \\
    \rowcolor{rowgreen}
    3 & Raw          & Corrected & Post-split \\
    4 & Preprocessed & Paper     & None       \\
    \rowcolor{roworange}
    5 & Preprocessed & Paper     & Pre-split  \\
    6 & Preprocessed & Corrected & None       \\
    7 & Preprocessed & Corrected & Pre-split  \\
    \bottomrule
  \end{tabular}
  \begin{tablenotes}
    \footnotesize
    \item[\textcolor{rowgreen!80!black}{$\blacksquare$}]
      Condition~3: fully corrected — raw data eliminates the
      resolution-confound amplification introduced by preprocessing.
    \item[\textcolor{roworange!80!black}{$\blacksquare$}]
      Condition~5: replication of Guha et al.\ (2025) exact setup
      (preprocessed, paper labels, pre-split downsampling).
  \end{tablenotes}
\end{threeparttable}
\end{table}
